\documentclass{article}
\usepackage{../math_template}

\author{Jonathan Kasongo}
\title{British Math Olympiad 2013 Round 1 --- P3/6}

\begin{document}
\maketitle

\begin{problem}{British Math Olympiad 2013 Round 1 --- P3/6}
A number written in base 10 is a string of $3^{2013}$ digit 3s. No other
digit appears. Find the highest power of 3 which divides this number.
\end{problem}

\begin{solution}{Solution}
Let $N = 3\cdot 10^{3^{2013} - 1} + 3\cdot 10^{3^{2013} - 2} +
\cdots + 3\cdot 10^1 + 3\cdot 10^0 = 3\cdot
\frac{10^{3^{2013}} - 1}{10 - 1}$. We claim that $v_3 (N) = 3^{2013}$.
It suffices to show that $v_3 \left(\frac{1}{3} \left( 10^{3^{2013}} - 1
\right) \right) = 3^{2013} \iff v_3 \left(10^{3^{2013}} - 1\right) =
{3^{2013} + 1}$. \\

First we show that $10^{3^{2013}} \equiv 1 \ \left(\text{mod} \
3^{3^{2013}+1}\right)$. Note that a positive integer is coprime with
$3^{3^{2013}+1}$ as long as it doesn't have a factor of 3. We know that
every third number of the
positive integers is a multiple of 3, and thus the number of positive
integers that are coprime with $3^{3^{2013} + 1}$ is exactly $\frac{2}{3} \cdot
3^{3^{2013} + 1} = 2 \cdot 3^{3^{2013}}$ and thus $\varphi \left(
3^{3^{2013}+1}\right) = 2 \cdot 3^{3^{2013}}$. Now notice that
$1 \equiv 10^{2\cdot 3^{2013}} \equiv 10^{3^{2013}} \ \left(\text{mod} \
3^{3^{2013}+1} \right)$ by Euler's theorem, and this establishes
that $3^{3^{2013} + 1} \mid 10^{3^{2013}} - 1$ as desired. \\

Now we just need to show that $3^{3^{2013} + 2} \nmid 10^{3^{2013}} - 1$.
By the same arguement as last time we know that
$\varphi \left( 3^{3^{2013} + 2} \right) = \frac{2}{3} \cdot
3^{3^{2013} + 2} = 2 \cdot 3^{3^{2013} + 1}$. Assume, for the sake of
contradiction, that $3^{3^{2013} + 2} \mid 10^{3^{2013}} - 1$. Thus
$10^{3^{2013}} - 1 \equiv 3^{2\cdot 3^{2013}} - 1 \equiv
3^{(3^{2013} + 2) + (3^{2013} - 2)} - 1 \equiv -1 \equiv 0 \
\left(\text{mod }  3^{3^{2013} + 2}\right)$ which is a contradiction.
Thus $3^{3^{2013} + 2} \nmid 10^{3^{2013}} - 1$ and so
$v_3 \left(10^{3^{2013}} - 1\right) = {3^{2013} + 1}$ as desired. Thus the
highest power of 3 dividing $N$ is $3^{2013}$. $\blacksquare$
\end{solution}

\end{document}
